% ============================================================================
%  1. INTRODUCCI\'ON
%  ----------------------------------------------------------------------------
%  Contextualiza el conjunto de las actividades a4 (correo y archivos),
%  a5 (DHCP y Web), a7 (Samba 4 AD DC) y a8 (DNS del dominio y Web integrado).
% ============================================================================

\section{Introducci\'on}

En toda organizaci\'on, la infraestructura de red constituye el pilar
sobre el cual se sustentan las operaciones diarias. Un banco, con
independencia de su tama\~no, depende de que sus empleados puedan resolver
los nombres de los servicios que utilizan, de que las estaciones obtengan
autom\'aticamente su configuraci\'on de red al conectarse y de que los
sistemas internos resulten inaccesibles desde el exterior sin una
autorizaci\'on expresa. Estos requisitos, que en el uso cotidiano pasan
desapercibidos, se traducen t\'ecnicamente en un conjunto de servicios
fundamentales que deben operar de forma coordinada y segura
\cite{tanenbaum2017}.

La resoluci\'on de nombres de dominio (\textsc{dns}) constituye la pieza
central de esa coordinaci\'on. El protocolo traduce los nombres que los
usuarios recuerdan, como \texttt{web.sudoers.lan}, en las direcciones
\texttt{ip} que las m\'aquinas entienden \cite{rfc1034, rfc1035}. Sobre esa
base se articulan el resto de los servicios: el correo electr\'onico
necesita localizar al servidor de entrega a trav\'es de registros
\textsc{mx} \cite{postel1982, criden2005}; la asignaci\'on din\'amica de
direcciones evita configurar cada estaci\'n a mano \cite{rfc2131, droms1997};
y la autenticaci\'on de los usuarios se apoya en un directorio que debe ser
localizada de forma inequ\'ivoca \cite{rfc4511, rfc4120}.

\subsection{La importancia de la administraci\'on de servidores GNU/Linux}

Las organizaciones actuales operan su infraestructura sobre sistemas
libres, y no por una cuesti\'on de preferencia tecnol\'ogica sino por
razones operativas y econ\'omicas bien documentadas. Un sistema GNU/Linux
ofrece un modelo de seguridad robusto basado en la separaci\'on de
privilegios, el control de acceso discrecional y la trazabilidad de las
operaciones, principios que se aplican de forma sistem\'atica en la
administraci\'on de servidores \cite{stewart2017}. Adem\'as, la
disponibilidad del c\'odigo fuente permite auditar el comportamiento de los
servicios criticales y corregir vulnerabilidades sin depender de los
plazos del proveedor.

Administrar estos servidores exige un conjunto amplio de competencias:
configurar servicios de red, gestionar usuarios y permisos, resolver
incidencias y aplicar medidas de securizaci\'on que reduzcan la
superficie de ataque \cite{nemeth2017}. El presente trabajo final
recoge justamente ese conjunto de competencias aplicado a un escenario
empresarial concreto.

\subsection{Antecedentes y contexto del proyecto}

El proyecto se desarrollo de manera incremental a lo largo de las
actividades del curso, y cada etapa aport\'o una pieza distinta a la
arquitectura final. La primera etapa implement\'o el servicio de correo
electr\'onico con \textsc{postfix} y \textsc{dovecot}, junto con un servidor
de archivos multi-protocolo que combina \textsc{smb}, \textsc{nfs} y
\textsc{http} sobre un mismo directorio compartido. La segunda etapa
incorpor\'o el servidor de asignaci\'on din\'amica de direcciones con
\textsc{kea} y un servidor web basado en \textsc{nginx} que act\'ua como
\textit{proxy inverso} \cite{keaDocs, nginxDocs}. La tercera etapa
despleg\'o el controlador de dominio \textsc{samba 4} \textsc{active
directory}, que centraliza la identidad, la autenticaci\'on Kerberos y la
compartici\'on de archivos \cite{sambateam2024, msft_activedirectory}. La
cuarta y \'ultima etapa consolid\'o el \textsc{dns} autoritativo del dominio
servido por el propio controlador y la publicaci\'on de los servicios por
nombre a trav\'es del proxy, con cifrado \textsc{https} y una autoridad de
certificaci\'on interna \cite{iscbind9, opensslDocs}.

\subsection{Problem\'atica}

La organizaci\'on se apoya en un \'unico servidor Debian que
concentra la identidad, los archivos, el correo y la publicaci\'on web. En
ese escenario se plantearon tres carencias que limitaban la operaci\'on:

\begin{itemize}
  \item \textbf{Resoluci\'on de nombres incompleta:} sin un servidor
        \textsc{dns} propio, cada cliente deb\'ia memorizar direcciones
        \textsc{ip} o depender de archivos \texttt{/etc/hosts} dispersos,
        esquema que no escala y que impide que un mismo nombre se resuelva
        de forma consistente en toda la red.
  \item \textbf{Asignaci\'on de direcciones manual:} configurar la red de
        cada estaci\'n de forma individual consume tiempo y es propenso a
        errores, adem\'s de impedir la gesti\'on centralizada de los
        par\'ametros de la red.
  \item \textbf{Ausencia de una puerta \'unica y segura:} sin un servidor web
        que concentre la publicaci\'on de los servicios, cada aplicaci\'on
        deber\'ia exponer su propio puerto, sin una capa de cifrado
        unificada ni un punto \'unico de control y de auditor\'ia.
\end{itemize}

\subsection{Prop\'osito del trabajo final}

El prop\'osito de este trabajo es demostrar la capacidad de dise\~nar,
implementar y validar una arquitectura de red empresarial completa sobre
GNU/Linux, integrando los servicios de directorio, asignaci\'on de direcciones,
resoluci\'on de nombres, publicaci\'on web, correo electr\'onico y
archivos, con medidas de securizaci\'on aplicadas en cada uno de ellos
\cite{debianReference, hertzog2020}. El objetivo no es simplemente que los
servicios funcionen, sino que operen de forma integrada: que un cliente
reciba su configuraci\'on por \textsc{dhcp}, resuelva los nombres contra el
controlador de dominio, se autentique mediante \textsc{kerberos} y
acceda a los servicios publicados por nombre a trav\'es de un \'unico
punto de entrada cifrado.

\subsection{Marco te\'orico}

\subsubsection{El sistema de nombres de dominio}

El \textsc{dns} es un servicio distribuido y jer\'arquico que traduce
nombres de dominio a direcciones \textsc{ip} \cite{rfc1034}. En un dominio
de \textsc{active directory} no es un accesorio, sino el transporte del
directorio: los clientes lo consultan para descubrir d\'onde reside el
servidor \textsc{ldap}, el \textsc{kerberos} o el propio controlador
mediante registros \textsc{srv} \cite{rfc2782}. Por esta raz\'on, un
controlador de dominio debe ser autoritativo de su zona
\cite{msft_activedirectory}. Los registros empleados en el proyecto son
los de tipo \textsc{a}, que asocian un nombre a una direcci\'on; los
\textsc{srv}, que indican qu\'e servidor atiende qu\'e servicio en qu\'e
puerto; y los registros \textsc{ns} y \textsc{soa}, que declaran la
autoridad y la vigencia de la zona.

\subsubsection{La asignaci\'on din\'amica de direcciones}

El protocolo \textsc{dhcp} permite que un servidor asigne direcciones
\textsc{ip} de forma autom\'atica a los clientes que se conectan
\cite{rfc2131}. El proceso sigue una secuencia de cuatro mensajes en la que
el cliente descubre, ofrece, solicita y confirma una direcci\'on. Adem\'as
de la direcci\'on, el servidor entrega par\'ametros de configuraci\'on
como la m\'ascara de subred, la puerta de enlace y el servidor
\textsc{dns}, lo que evita la intervenci\'on manual \cite{droms1997}. El
servidor \textsc{kea} se escogi\'o por su soporte activo y su
configuraci\'on basada en \textsc{json} \cite{keaDocs}.

\subsubsection{El servidor web y el proxy inverso}

Un servidor web atiende peticiones \textsc{http} y entrega documentos a
los navegadores. En el proyecto se emplea \textsc{nginx} con una doble
funci\'on: servir el contenido del portal y actuar como \textit{proxy
inverso}, es decir, recibir las peticiones y reenviarlas al servicio
interno correspondiente seg\'un el nombre solicitado \cite{nginxDocs}. El
proxy termina adem\'as la conexi\'on \textsc{tls}, de modo que el
cifrado se aplica en el punto \'unico de entrada mientras que las
comunicaciones internas se realizan en claro.

\subsubsection{La identidad centralizada}

El \textsc{active directory} centraliza la gesti\'on de usuarios, grupos y
permisos sobre la base de \textsc{ldap} \cite{rfc4511}, y delega la
autenticaci\'on en \textsc{kerberos}, cuyo protocolo de tickets evita
transmitir la clave del usuario en cada petici\'on \cite{rfc4120}. Esta
arquitectura es la que permite, en lugar de credenciales dispersas por
cada servicio, una gesti\'on \'unica y coherente de la seguridad
\cite{sambateam2024}.

\subsection{Estructura del documento}

El presente informe se organiza en seis cap\'itulos. El cap\'itulo
siguiente enuncia los objetivos del trabajo y su alcance. El tercero
documenta el diagrama arquitect\'onico final, con la segmentaci\'on de la
red, los roles de cada servidor y la justificaci\'on del dise\~no. El
cuarto desarrolla t\'ecnicamente cada servicio desplegado, siguiendo el
orden en que se implementaron. El quinto expone los desaf\'ios
encontrados durante la ejecuci\'on y las soluciones aplicadas. El sexto
expone las conclusiones, los aprendizajes y las recomendaciones para
futuras implementaciones.
